
\begin{longtable}[c]{ |p{1cm}|p{8cm}|p{1cm}|p{1cm}|p{1cm}|p{1cm}|}
    \caption{死区生成器的典型操作模式}
    \label{table:mod_dt}
    \\\hline
    \rowcolor{lightgray}
   模式
    & 描述
    & S0
    & S1
    & S2
    & S3
    \\\hline
    \endfirsthead
    \hline
    \rowcolor{lightgray}   
   模式
    & 描述
    & S0
    & S1
    & S2
    & S3
    \\\hline
    \endhead
    \hline
    \endfoot
    1
    & PWM\textcolor{red}{\it{x}}A 和 PWM\textcolor{red}{\it{x}}B 波形无变化
    & 1
    & 1
    & X
    & X
    \\\hline
    2
    & 高电平有效互补 (AHC)，参见图 \ref{figure:wf_dt_ahc}
    & 0
    & 0
    & 0
    & 1
    \\\hline
    3
    & 低电平有效互补 (ALC)，参见图 \ref{figure:wf_dt_alc}
    & 0
    & 0
    & 1
    & 0
    \\\hline
    4
    & 高电平有效 (AH)，参见图 \ref{figure:wf_dt_ah}
    & 0
    & 0
    & 0
    & 0
    \\\hline
    5
    & 低电平有效 (AL)，参见图 \ref{figure:wf_dt_al}
    & 0
    & 0
    & 1
    & 1
    \\\hline
    6
    & PWM\textcolor{red}{\it{x}}A 输出 = PWM\textcolor{red}{\it{x}}A 输入（无延迟）
    & 0
    & 1
    & 0 或 1
    & 0 或 1 \\
    & PWM\textcolor{red}{\it{x}}B 输出 = PWM\textcolor{red}{\it{x}}A 输入，下降沿延迟&&&&
    \\\hline
    7
    & PWM\textcolor{red}{\it{x}}A 输出 = PWM\textcolor{red}{\it{x}}A 输入，上升沿延迟
    & 1
    & 0
    & 0 或 1
    & 0 或 1\\
    & PWM\textcolor{red}{\it{x}}B 输出 = PWM\textcolor{red}{\it{x}}B 输入（无延迟） &&&&
    \\\hline
    \multicolumn{6}{|l|}{说明：以上所有模式中，S4 - S8 的开关位置都置 0。}
\end{longtable}